% =============================================================================
%  Sección 5: Gestión
% =============================================================================
\section{Fase 5: Viabilidad y Gestión}
\label{sec:viabilidad-gestion}

La fase final de este Plan Maestro define los mecanismos operativos, jurídicos 
y financieros necesarios para transitar de la propuesta a la ejecución. 
En el contexto de Milpa Alta, la gestión se enfrenta a un desafío 
estructural: la \textbf{marginalización geográfica} y política de una 
periferia que históricamente ha sido tratada con una ruralidad desatendida 
por el gobierno central de la Ciudad de México \cite{milpaalta2017pddu}.

\subsection{El Reto de la Descentralización y los Recursos}
\label{ssec:reto-descentralizacion}

San Pedro Atocpan, aunque es el ejemplo más visible por su potencia comercial, 
comparte con los otros 11 pueblos de la alcaldía una vulnerabilidad derivada de 
la lejanía respecto al centro político de la capital. 

\begin{itemize}
    \item \textbf{Inequidad Presupuestal:} Las políticas de inversión en la 
        Ciudad de México han tendido históricamente a favorecer 
        \termino{proyectos megaespeciales} en zonas centrales o de alta 
        plusvalía, relegando a las alcaldías rurales a presupuestos 
        insuficientes para la conservación de su vasto 
        patrimonio \cite{balandrano2016conservacion}. 
    \item \textbf{Debilidad Institucional:} La gestión local padece 
        una crónica escasez de recursos propios y personal especializado, 
        lo que, sumado a la falta de coordinación entre las oficinas de 
        conservación y las autoridades locales, impide un embellecimiento 
        urbano continuo y coherente en la demarcación \cite{balandrano2016conservacion}.
\end{itemize}

\subsection{Alineación Normativa y Gobernanza Comunal}
\label{ssec:alineacion-normativa-gobernanza}

Para superar el \textbf{"nudo ciego jurídico"} que detiene el desarrollo en la zona, 
la viabilidad del plan descansa en un modelo de 
\textbf{Crecimiento Endógeno Controlado}.

\begin{enumerate}
    \item \textbf{Certeza de Posesión:} Es imperativo que la Asamblea Comunal de 
            San Pedro Atocpan valide la utilidad pública de las 
            intervenciones (como el CETRAM o la infraestructura hídrica) 
            para otorgar certeza jurídica sin vulnerar el régimen 
            inalienable de la tierra \cite{milpaalta2017pddu}.
    \item \textbf{Instrumentos de Gestión:} Se propone la creación de 
            un \textit{Consejo de Gestión de los 12 Pueblos}, que funcione 
            como un ente facilitador y articulador de intereses, evitando 
            que las decisiones se tomen de forma arbitraria o sesgada 
            desde la administración central \cite{balandrano2016conservacion}.
\end{enumerate}

\subsection{Ruta Crítica y Sostenibilidad Financiera}
\label{ssec:ruta-critica}

Debido a que el centralismo suele ir de la mano con la discontinuidad 
política (administraciones de 3 años), el plan propone etapas que 
trasciendan los tiempos electorales para evitar el abandono de 
programas válidos \cite{balandrano2016conservacion}.

\subsection{Estrategia Antigentrificación y Protección del Tejido Social}
\label{ssec:estrategia-proteccion-tejido-social}

Un riesgo latente en la rehabilitación de centros históricos es la expulsión 
de los habitantes originales por el aumento del valor 
del suelo \cite{balandrano2016conservacion}. En San Pedro Atocpan, 
la gestión debe priorizar:

\begin{itemize}
    \item El reconocimiento del \textbf{Uso Real} habitacional para 
            regularizar servicios básicos sin incentivar la especulación 
            inmobiliaria externa.
    \item El fortalecimiento de la producción de mole como motor económico 
            local, garantizando que las ganancias de la revitalización urbana 
            permanezcan en manos de las familias del pueblo originario.
\end{itemize}